\section{Conclusion}
\label{sec:conclusion}

We asked whether layer births come in different kinds that can be told apart by measurement. Within the Six-Birds framework the answer supported here is yes: at least four. All four classes form structure (packaging). They differ in whether drive is present at birth and whether the birth moves with the observation timescale (staging).

Each class has an explicit finite representative, and its label is certified exactly on the declared grids. The undriven staging class is certified at every admissible size. The structure-with-drive versus structure-without-drive contrast (Class-I versus Class-II) is the most robust result: it rests on an exactly zero versus a clearly positive affinity, and controls show that the drive channel does not create false arrows. The staging classes are confirmed under the manual lens only.

At a common size the four classes occupy distinct regions of a three-coordinate observable map built from the same quantities that define them. Capacity/depth proxies turn out to be a rescaling of the structural boundary and are not classifiers. Kernels built from upstream compiled features can be placed on the map and land in Class-I.

The result is a vocabulary and a set of measurements, not a theory of universality classes. Whether these four classes remain distinct in the asymptotic, exponent-level sense is the next question, and this paper does not answer it.
